%============================================================
% CH1 — SOLUTION SKETCHES FOR THE PRACTICE SET
% Idea, invariant or key fact, algorithm, complexity and (for the traps) the trap.
% Sketches, not full code: the point is to check your DESIGN.
% Every sketch was implemented in C++14 and reproduces the official samples;
% every non-trivial one was also compared with a brute force on random small inputs.
% Order and groups follow problem_set.tex (Tiers 1-3, then Stretch).
%============================================================
\subsection{Solution Sketches}
\label{ssec:ch1-sketches}

\begin{warningbox}[Read a sketch only after a real attempt]
Each sketch gives the \emph{design} --- the key fact, what the loop keeps, the algorithm, the complexity and, where there is one, the trap --- in a few lines, so you can compare it with your own. Writing and debugging the code is left to you. If your design differs but you can argue why it computes the answer, it is probably also right.
\end{warningbox}

\begingroup
\providecommand{\sk}{}
\renewcommand{\sk}[2]{\par\smallskip\noindent\textbf{\textcolor{deepblue}{#1}}\quad #2\par}

\begin{solutionbox}[Tier 1 --- Warm-up]
\sk{P1.01 Next Round (CF 158A).}{The scores are given in non-increasing order, so the $k$-th place score is simply $a_k$. Count the $i$ with $a_i \ge a_k$ \emph{and} $a_i > 0$. Ties with $a_k$ advance too, which is why the answer can exceed $k$ (sample~1: $k = 5$, answer $6$); and if $a_k = 0$, the zeros stay at home (sample~2: answer $0$). One pass, \bigO{n}.}
\sk{P1.02 Monotonic Array.}{Two flags, both \texttt{true} at the start: a step $a_i < a_{i-1}$ kills ``non-decreasing'', a step $a_i > a_{i-1}$ kills ``non-increasing''. Equal neighbours kill neither. Answer: at least one flag survives. One pass, \bigO{n}, no need to decide the direction in advance. Samples: \texttt{true}, \texttt{true}, \texttt{false}.}
\sk{P1.03 Sign of the Product of an Array.}{The state is one number $s \in \{-1, 0, 1\}$: start with $1$, a zero sets it to $0$ for good, a negative value flips it. \emph{Trap:} multiplying the values overflows after a handful of elements (and \texttt{double} loses the sign of an underflowed $0$ just as badly); the answer needs only the parity of the negatives and the presence of a zero. \bigO{n}. Samples: $1$, $0$, $-1$.}
\sk{P1.04 Stones on the Table (CF 266A).}{Imagine the compaction of Pattern~2 that keeps a stone only if it differs from the last \emph{kept} stone. In a run of equal colours, the first stone is kept and every other is removed, so the number of removals is the number of $i$ with $s_i = s_{i-1}$. Nothing needs to be moved. \bigO{n}. Samples: $1$, $4$, $0$.}
\sk{P1.05 Reverse Only Letters.}{The reversal loop with $l = 0$, $r = n-1$: if $s_l$ is not a letter, $l{+}{+}$; else if $s_r$ is not a letter, $r{-}{-}$; else swap and move both. Invariant: outside $[l, r]$ every letter is already in its final place and every other character never moved. \bigO{n}, in place. All three samples reproduced (\texttt{"dc-ba"}, \dots).}
\sk{P1.06 Normal Problem (CF 2044B).}{From inside the store the window is read right to left \emph{and} every letter is mirrored: \texttt{p} and \texttt{q} turn into each other, \texttt{w} stays. Reverse the string, then map \texttt{p}$\leftrightarrow$\texttt{q}. \bigO{n}. Sample: \texttt{qwq} $\to$ \texttt{pwp}, \texttt{ppppp} $\to$ \texttt{qqqqq}, \dots}
\sk{P1.07 Who's Opposite? (CF 1560B).}{On a circle of $n$ people (even $n$), person $x$ looks at $x + n/2$ or $x - n/2$, so $|a - b| = n/2$ and $n = 2|a-b|$ is forced. The opposite of $c$ is $c + n/2$ if $c \le n/2$, else $c - n/2$. \emph{Trap:} the circle must contain $a$, $b$ and $c$; if any of them exceeds $n$, print $-1$ (sample: $a=2$, $b=3$ gives $n=2$, which has no person $3$). \bigO{1} per test. All seven sample answers reproduced; checked against trying every even $n < 200$.}
\sk{P1.08 Maximum Difference Between Adjacent Elements in a Circular Array.}{Answer $\max_i |a_i - a_{(i+1) \bmod n}|$. The index formula makes the pair (last, first) one more ordinary pair --- no extra line after the loop, no special case for $n = 2$. \bigO{n}. Samples: $3$, $5$.}
\sk{P1.09 Helpful Maths (CF 339A).}{Counting sort on the value range $\{1,2,3\}$: \texttt{cnt[d]++} for every digit, then write \texttt{cnt[1]} ones, \texttt{cnt[2]} twos and \texttt{cnt[3]} threes, separated by \texttt{+}. \bigO{n} instead of \bigO{n \log n}, and the separators are produced while writing (no \texttt{+} before the first term). Samples: \texttt{1+2+3}, \texttt{1+1+1+3+3}, \texttt{2}.}
\sk{P1.10 Submission Bait (CF 1990A).}{Alice wins iff some value occurs an odd number of times. If every frequency is even, Bob copies Alice's move (same value) and Alice runs out first. Otherwise Alice takes the largest value $v$ with odd frequency: every larger value has even frequency, the rest of $v$ is even, and now she is the one who copies. \emph{Trap:} ``take the maximum'' loses on $[1, 2, 2]$ (Bob answers with the other $2$), while taking the $1$ first wins. A frequency array over $1..n$, \bigO{n}. Samples: YES NO YES NO YES; checked against an exhaustive game search for $n \le 6$.}
\end{solutionbox}

\begin{solutionbox}[Tier 2 --- Core]
\sk{P1.11 Good Prefixes (CF 1985C).}{If $x = S - x$ for some element $x$ of a prefix with sum $S$, then $x = S/2 \ge$ every other element (they are non-negative and sum to $x$), so $x$ is the maximum. Hence a prefix is good iff $2 \cdot \max = S$. Keep the running sum and the running maximum: \bigO{n}. Use \texttt{long long}: $S$ reaches $2 \cdot 10^5 \cdot 10^9$. Samples: $1, 0, 3, 3, 4, 1, 2$; checked against testing every element of every prefix.}
\sk{P1.12 Even-Odd Increments (CF 1744B).}{State: the sum $S$, the number $E$ of even and $O$ of odd elements. Query ``$0\ x$'': $S \mathrel{+}= E \cdot x$, and if $x$ is odd all evens become odd ($O \mathrel{+}= E$, $E = 0$); query ``$1\ x$'' symmetrically. The array itself is never updated: \bigO{1} per query, \bigO{n+q} in total, versus \bigO{nq} for the honest simulation. $S$ exceeds $2^{31}$ (sample: $3000000094$): \texttt{long long}. Checked against the simulation.}
\sk{P1.13 Remove One Element to Make the Array Strictly Increasing.}{Scan for $a_i \le a_{i-1}$. A second violation means \texttt{false}. At the first one, remove $a_{i-1}$ if $i < 2$ or $a_{i-2} < a_i$ (then $a_i$ stays as the last kept value); otherwise remove $a_i$, i.e.\ overwrite $a_i \gets a_{i-1}$ so that later comparisons see the kept value. \emph{Trap:} always deleting the left (or always the right) element fails --- $[1,2,10,5,7]$ needs the $10$ gone, $[2,3,1,4]$ needs the $1$ gone. \bigO{n}. Samples: \texttt{true}, \texttt{false}, \texttt{false}; checked against trying every deletion.}
\sk{P1.14 Minimum Deletions to Make Array Beautiful.}{Scan $i = 0..n-1$ and count deletions $d$. The kept array has $i - d$ elements before $a_i$, so $a_i$ would land at an even index of the kept array iff $i - d$ is even. If so and $a_i = a_{i+1}$, delete $a_i$ ($d{+}{+}$): deleting now is never worse, since the conflict at this position must be fixed anyway and later positions only shift. At the end, if $n - d$ is odd, delete the last element. \bigO{n}, \bigO{1} memory. Samples: $1$, $2$; checked against all subsequences for $n \le 9$.}
\sk{P1.15 Reverse (CF 1638A).}{Let $i$ be the first position with $p_i \ne i$ (if none, the identity is already optimal: reverse $[1,1]$). Positions before $i$ cannot improve, and the smallest value that can reach position $i$ is $i$ itself, from its position $j > i$. Reversing $[i, j]$ puts it there; among all reversals that do, this is the only one (the left end must be $i$, the right end must be $j$). \bigO{n}. Samples reproduced; checked against trying all $\binom{n+1}{2}$ reversals.}
\sk{P1.16 Entertainment in MAC (CF 1935A).}{Two reversals cancel, so with an even $n$ only two results matter: $s$ itself (reverse $n$ times) and $\mathrm{rev}(s) + s$ (one reversal, one append of the reversed string, then cancelling pairs). Any further append only makes the string longer with the same beginning, which is never lexicographically smaller. Answer $\min(s, \mathrm{rev}(s) + s)$. \emph{Trap:} ``a shorter string is smaller'' is wrong --- for $s = \texttt{grib}$ the answer is the longer \texttt{birggrib}, because it starts with $\mathrm{rev}(s) = \texttt{birg} < \texttt{grib}$ (sample~2). Appending to $s$ itself never helps: $s$ is a prefix of $s + \mathrm{rev}(s)$. \bigO{|s|}. All five samples reproduced; checked against all $2^n$ operation sequences for $n \le 6$.}
\sk{P1.17 Symmetric Encoding (CF 1974B).}{A presence table \texttt{has[26]} gives the distinct letters $r_0 < r_1 < \dots < r_{m-1}$ in alphabetical order without sorting. The encoding maps $r_k \mapsto r_{m-1-k}$; the encoded string has the same set of letters, and the map is its own inverse (mirroring twice is the identity), so decoding applies the same map. \bigO{n + 26}. Samples: \texttt{serofedsoc} $\to$ \texttt{codeforces}, \dots; checked by encoding random strings and decoding them back.}
\sk{P1.18 Can Place Flowers.}{Scan left to right; plant at $i$ when $f_i = 0$ and both neighbours are empty, treating a missing neighbour (at $i = 0$ or $i = m-1$) as empty, and write the $1$ into the bed immediately so the next position sees it. Planting as early as possible never blocks more positions than planting later. \bigO{m}. Samples: \texttt{true}, \texttt{false}; checked against all subsets of positions for $m \le 9$.}
\sk{P1.19 Two Friends (CF 1969A).}{Inviting $i$, $p_i$ and $p_{p_i}$ always brings $i$ and $p_i$: so the answer is at most $3$. Two invitations $\{x, y\}$ bring two people only if $p_x = y$ and $p_y = x$. So the answer is $2$ iff some $i$ has $p_{p_i} = i$ (a double index lookup, as in Problem~\ref{ssec:ch1-build-permutation}), else $3$. \bigO{n}. Samples: $2, 3, 2$; checked against all subsets for $n \le 8$.}
\sk{P1.20 Yet Another Palindrome Problem (CF 1324B).}{A length-$3$ palindromic subsequence is $x\,y\,x$: two equal values at distance $\ge 2$. Record the first occurrence \texttt{first[v]} of each value; at position $j$, if \texttt{first[}$a_j$\texttt{]} $\le j - 2$, answer YES. Keeping the \emph{first} occurrence is enough: it maximises the distance. \emph{Trap:} equal neighbours are not enough --- $[1, 1, 2]$ is NO, while $[1, 2, 2, 1]$ is YES through the two $1$s. \bigO{n}. Samples: YES YES NO YES NO; checked against all triples.}
\sk{P1.21 Remove Prefix (CF 1714B).}{The array left after removing $k$ elements is the suffix $a[k..n)$, which must have distinct values; the shortest removal is the longest distinct suffix. Scan from the back with a \texttt{seen} array over $1..n$, stop at the first value already seen at index $i$; answer $i + 1$ (or $0$ if the scan finishes). \bigO{n}. Samples: $1, 4, 0, 0, 5$; checked against removing one element at a time.}
\end{solutionbox}

\begin{solutionbox}[Tier 3 --- Challenge]
\sk{P1.22 Make It Good (CF 1385C).}{Taking the smaller end each time produces a non-decreasing sequence iff the array rises to a maximum and then falls (a mountain): the smaller end is always the next value of the sorted order. So the answer is $n$ minus the longest mountain \emph{suffix}. From $i = n-1$: move left while $a_{i-1} \ge a_i$, then while $a_{i-1} \le a_i$; print $i$. \bigO{n}. Samples: $0, 4, 0, 2, 3$; checked against simulating the removal process on every suffix.}
\sk{P1.23 Find the Winner of an Array Game.}{The front element is always the winner of all previous rounds, so the rounds just scan the array left to right: keep the champion and its streak; a new larger element becomes champion with streak $1$, otherwise the streak grows. Stop when the streak reaches $k$. If the scan ends first, the champion is the maximum, which wins every later round. \emph{Trap:} $k$ up to $10^9$ --- simulating the rotating queue is hopeless and unnecessary. \bigO{n}. Samples: $5$, $3$; checked against a direct queue simulation.}
\sk{P1.24 Maximum Number of Operations to Move Ones to the End.}{An operation moves a one across one whole block of zeros. To maximise the number of operations, move the leftmost movable ones first, so that ones pile up and cross every later block one by one: each block of zeros is then crossed once by every one to its left. Answer $\sum$ over zero blocks of (ones before the block): add the running count of ones whenever $s_i = 0$ and $s_{i-1} = 1$. \bigO{n}, \texttt{long long} ($10^5$ ones times $5 \cdot 10^4$ blocks). Samples: $4$, $0$; checked against an exhaustive search over all operation orders.}
\sk{P1.25 Array Fix (CF 1948B).}{Scan left to right keeping \texttt{prev}, the last value written to the result. For $a_i \ge 10$ with tens digit $\le$ ones digit and tens digit $\ge$ \texttt{prev}, split it (\texttt{prev} $\gets$ ones digit $\le a_i$); otherwise keep $a_i$ if $a_i \ge$ \texttt{prev}; otherwise answer NO. Splitting whenever allowed is safe because it leaves a smaller \texttt{prev}, which can only help later elements. \bigO{n}. Samples: YES NO YES; checked against all $2^n$ split choices.}
\sk{P1.26 Wrong Addition (CF 1619C).}{Walk from the right of $a$ and $s$. For digit $x$ of $a$, take the last unused digit $y$ of $s$; if $y \ge x$, the digit of $b$ is $y - x$; else take one more digit of $s$ to form $10 \le y \le 19$ and require $y - x \le 9$. When $a$ is exhausted, the remaining digits of $s$ are the leading digits of $b$; strip leading zeros. \emph{Trap:} answer $-1$ when $s$ runs out of digits or a two-digit $y - x$ exceeds $9$ --- do not assume an answer exists, and read $a$, $s$ as strings (up to $10^{18}$). \bigO{\text{digits}}. Samples: $3465, 4, -1, 90007, 10, -1$; checked against trying every $b < 2 \cdot 10^4$.}
\sk{P1.27 Deja Vu (CF 1891B).}{After query $x$, an element it touched has the form $m 2^x + 2^{x-1}$: divisible by $2^{x-1}$ but not by $2^x$, so no query $x' \ge x$ ever touches it again; untouched elements were already not divisible by $2^x$. Therefore only queries strictly smaller than every earlier query matter, at most $30$ of them; apply each to the whole array. \emph{Trap:} the naive \bigO{nq} simulation is $10^{10}$. \bigO{30n + q}. All four samples reproduced; checked against the naive simulation.}
\sk{P1.28 Pleasant Pairs (CF 1541B).}{$a_i a_j = i + j$ forces $i + j \equiv 0 \pmod{a_i}$. For each $i$, start at the smallest $j > i$ with $(i + j) \bmod a_i = 0$ and step $j \mathrel{+}= a_i$ up to $n$, testing $a_i a_j = i + j$. Since the values are distinct, the number of steps is at most $\sum_i n / a_i \le n \cdot H_{2n} = \bigO{n \log n}$. Use \texttt{long long} for the product. Samples: $1, 1, 3$; checked against all pairs.}
\sk{P1.29 Flower City Fence (CF 1862C).}{Laid horizontally, plank row $j$ has length $c_j = \#\{i : a_i \ge j\}$; the fence is symmetric iff $c_j = a_j$ for all $j$ and there are exactly $n$ rows, i.e.\ $a_1 = n$. Check $a_1 = n$ first (this also makes every $a_i \le n$, so a frequency array of size $n+1$ suffices), then count by value and accumulate from $j = n$ down: $c_j = c_{j+1} + \texttt{cnt}[j]$. \bigO{n}. Samples: YES YES NO NO YES YES NO; checked against building the horizontal fence.}
\end{solutionbox}

\begin{solutionbox}[Stretch --- above this chapter]
\sk{P1.30 Rotation Matching (CF 1365C).}{Rotating only $a$ is enough (rotating $b$ left equals rotating $a$ right). Value $v$ matches after the rotation that moves position $\mathrm{pos}_a[v]$ onto $\mathrm{pos}_b[v]$: shift $(\mathrm{pos}_a[v] - \mathrm{pos}_b[v]) \bmod n$. Count the shifts of all $n$ values in a frequency array and print the largest count. \bigO{n}. Samples: $5$, $1$, $2$; checked against trying all pairs of rotations.}
\sk{P1.31 Mad MAD Sum (CF 1990C).}{One round, done with a \texttt{seen} array, is \bigO{n}. After the first round the array is non-decreasing; after the second, every non-zero value appears at least twice in a row (except possibly at the end). From then on each round just shifts the array right by one cell (dropping the last), so element $b_i$ (0-indexed) is added in exactly $n - i$ more rounds. Answer: the sums of the first two rounds plus $\sum_i b_i (n - i)$, in \texttt{long long}. \bigO{n}. Samples: $1, 13, 9, 40$; checked against the full simulation.}
\sk{P1.32 Cobb (CF 1554B).}{$a_i | a_j < 2n$, so the penalty is below $2kn$ and $f(n-1, n) > n(n-1) - 2kn$. Any pair with $i < n - 2k$ has $f \le i j \le (n - 2k - 1) n < n(n-1) - 2kn$. So only $i, j \in [n-2k, n]$ need to be tried: $\bigO{k^2}$ pairs per test. \emph{Trap:} the \bigO{n^2} loop over all pairs is $10^{10}$; the bound comes from the indices, not the values. Samples: $-1, -4, 3, 12$; checked against all pairs.}
\end{solutionbox}

%------------------------------------------------------------
\subsection*{Answers to the Self-Check Quiz}
\label{ssec:ch1-quiz-answers}

\begin{solutionbox}[Quiz answers]
\begin{enumerate}[leftmargin=*]
    \item The elements are stored contiguously and all have the same size, so the address of $a[i]$ is $\text{base} + i \cdot \text{size}$: one multiplication and one addition, whatever $i$ is. A linked list has no such formula; it must follow $i$ pointers.
    \item \bigO{n} in total, \bigO{1} amortised per call. When the capacity is exceeded the vector doubles it and copies everything, but the copies cost $1 + 2 + 4 + \dots < 2n$ over $n$ pushes. Calling \texttt{reserve(n)} first removes the copies altogether.
    \item $[l, r)$ has $r - l$ elements; it is empty when $l = r$; it splits at any $m$ with $l \le m \le r$ into $[l, m)$ and $[m, r)$ with no overlap and no gap. The loop is \texttt{for (int i = l; i < r; ++i)}.
    \item About $10^8$ simple operations per second: \bigO{n} and \bigO{n \log n} are fine, \bigO{n^2} $= 4 \cdot 10^{10}$ is not. The sum can reach $2 \cdot 10^{14} > 2^{31}$, so it needs \texttt{long long} (the individual values still fit in \texttt{int}).
    \item $n(n+1)/2$ non-empty subarrays (choose the two ends), $2^n$ subsequences including the empty one (each element in or out). For $n = 2 \cdot 10^5$ that is about $2 \cdot 10^{10}$ subarrays --- so a solution may not even list them.
    \item A value equal to one already stored must be \emph{ignored}: otherwise $[5, 5, 4]$ would store $m_1 = m_2 = 5$ and report the wrong third maximum. The comparisons for equality come first, then the shifting down (\S\ref{ssec:ch1-p1-theory}).
    \item The written prefix $a[0, w)$ holds exactly the non-zero values among $a[0, r)$, in their original order, and $w \le r$. Since the write index never passes the read index, no value is overwritten before it is read; after the scan, $a[w, n)$ is filled with zeros.
    \item Reverse the whole array, then the first $k$ cells, then the last $n - k$ (for a rotation to the right). First replace $k$ by $k \bmod n$: rotating by $n$ changes nothing, and $k > n$ would make the second reversal run out of the array.
    \item Count the descents $a_i > a_{(i+1) \bmod n}$ \emph{including} the pair (last, first). The array is a rotation of a sorted array iff there is at most one. Forgetting the circular pair accepts wrong arrays: $[3, 4, 1, 2, 5]$ has one linear descent ($4 > 1$) but a second one from $5 > 3$, and it is not a rotation of $[1, 2, 3, 4, 5]$.
    \item Only when every digit is $9$: the carry survives the whole loop. Then the answer is $1$ followed by $n$ zeros; in every other case the loop stops at the first digit that is not $9$.
    \item Every cell must hold its old value and its new value at the same time. With all values in $[0, n)$, store $a[i] + n \cdot \text{new}_i$: the old value is $a[i] \bmod n$, the new one $\lfloor a[i] / n \rfloor$. The \texttt{\% n} is needed when reading $a[a[i]]$ because that cell may already have been encoded. The largest stored value is below $n^2$, so it must fit in the integer type.
    \item When the values lie in a small range $[0, V]$ with $V$ about $10^6$--$10^7$ at most: an array \texttt{cnt[V+1]} is faster and simpler than any map. For values up to $10^9$, use a hash map (Chapter~\ref{ch:hashing}), sorting (Chapter~\ref{ch:sorting}) or coordinate compression; an array of $10^9$ counters does not fit in memory.
\end{enumerate}
\end{solutionbox}
\endgroup
